\label{sec:method}

\paragraph{Problem setup.}
We consider trip-scale route generation on a fixed urban region.
Given an origin $o\in\mathbb{R}^2$, destination $d\in\mathbb{R}^2$, start time $t$ (Unix epoch seconds), and semantic context $s$ (e.g., POI density, road hierarchy, and other raster channels), the goal is to generate a trajectory
$x_{1:F}\in(\mathbb{R}^2)^F$ that is both \emph{structurally realistic} (on-road, smooth) and \emph{behaviorally plausible} (captures multi-modal corridor choice conditioned on $s$ and $t$).

\paragraph{Behavioral realism as context-conditioned distributions.}
Let $\mathcal{C}$ be a corridor set for a given OD constructed from the road graph.
We treat behavioral realism as matching the \emph{corridor choice distribution} under context:
different $(t,s)$ should induce different $p(z\mid o,d,t,s)$, and generated samples should reproduce the GT distribution shift across contexts (Sec.~\ref{sec:experiments}).
This explicitly distinguishes meaningful diversity from indiscriminate randomness: a model that produces diverse routes without conditioning on context is not behaviorally realistic---it is merely stochastic.

\paragraph{Road topology as a discrete corridor space.}
Let $G=(V,E)$ be a road graph extracted from OpenStreetMap.
We map $o$ and $d$ to nearby nodes (or projected points on edges) and compute a set of $K$ candidate corridors $\mathcal{C}=\{c_1,\dots,c_K\}$ using $k$-shortest paths under a travel cost (e.g., length or road-class weighted length).
Each corridor $c_k$ is represented by its polyline geometry and/or a sequence of edges.

\paragraph{CascadeTraj: decision then execution.}
CascadeTraj decomposes generation into two stages:
a decision stage that selects a corridor index $z\in\{1,\dots,K\}$,
and an execution stage that generates a continuous trajectory conditioned on the chosen corridor.
Formally,
\[
  p(x_{1:F}\mid o,d,t,s)\;=\;\sum_{z=1}^K p(z\mid o,d,t,s)\;p(x_{1:F}\mid z,o,d,t,s).
\]

\paragraph{Decision stage: corridor selection on the road graph.}
The decision model predicts $p(z\mid o,d,t,s)$ by scoring each candidate corridor with temporal--spatial semantics and OD context.
Semantics enter both \emph{globally} (pooling raster channels around the OD box) and \emph{corridor-wise} (pooling semantic channels along each corridor polyline).
This is where temporal--spatial semantics act as first-class conditioning signals: the same OD under different $(t,s)$ should produce different corridor distributions, and training aligns predicted distributions with GT corridor assignments derived from trajectory--corridor matching.

\paragraph{Execution stage: trajectory generation within a corridor.}
Conditioned on a corridor polyline, execution generates a dense trajectory using a residual diffusion model.
The corridor provides a coarse skeleton; the diffusion residual models within-corridor variations (lane choice, local detours, smoothing).
This stage is compatible with additional semantic/temporal conditioning and avoids post-hoc snapping that would distort corridor choice distributions.

\begin{figure}[t]
  \centering
  \maybeincludegraphics[width=\linewidth]{figures/fig3_method_comparison.pdf}
  \caption{\textbf{Method overview.} CascadeTraj separates corridor selection on the road graph from continuous trajectory execution, with temporal--spatial semantics conditioning the corridor decision.}
  \label{fig:method-comparison}
\end{figure}

\begin{figure}[t]
  \centering
  \maybeincludegraphics[width=\linewidth]{figures/fig4_road_graph.pdf}
  \caption{\textbf{Road graph corridor space.} Detroit road network with an example OD pair and $K$ candidate corridors (e.g., $k$-shortest paths).}
  \label{fig:road-graph}
\end{figure}
